% =========================================================================
% บทที่ 6: กลศาสตร์ควอนตัม แบบจำลองอะตอม และทฤษฎีการกระเจิง
% =========================================================================

\chapter{กลศาสตร์ควอนตัม แบบจำลองอะตอม และทฤษฎีการกระเจิง}
\label{ch:chapter06}

\noindent{\large\bfseries\color{physicsblue} สมการชเรอดิงเงอร์ที่ไม่ขึ้นกับเวลา บ่อศักย์อนันต์ และปรากฏการณ์ทะลุผ่านกำแพงศักย์}

\vspace{0.4cm}

\begin{equationsbox}{สมการชเรอดิงเงอร์และเงื่อนไขความน่าจะเป็น}
\begin{align}
-\frac{\hbar^2}{2m}\frac{d^2\psi(x)}{dx^2} + V(x)\psi(x) &= E\psi(x), \quad \int_{-\infty}^{\infty} |\psi(x)|^2 dx = 1 \\
T &\approx e^{-2\kappa L}, \quad \kappa = \frac{\sqrt{2m(V_0 - E)}}{\hbar}
\end{align}
\end{equationsbox}

\section{ทฤษฎีและกรอบแนวคิดสำคัญ}

กลศาสตร์ควอนตัมปฏิวัติแนวคิดดั้งเดิมเรื่องวิถีการเคลื่อนที่ของอนุภาค โดยแทนที่สถานะของระบบด้วยฟังก์ชันคลื่น (Wavefunction: $\psi$) ซึ่งความหนาแน่นความน่าจะเป็น (Probability Density) ในการพบอนุภาคถูกกำหนดโดยกฎของบอร์น (Born's Rule: $P(x) = |\psi(x)|^2$)

หนึ่งในปรากฏการณ์เชิงควอนตัมที่ไม่มีในกลศาสตร์คลาสสิกคือ การทะลุผ่านกำแพงศักย์ (Quantum Tunneling) ซึ่งอนุภาคที่มีพลังงาน $E$ น้อยกว่าความสูงของกำแพงศักย์ $V_0$ ยังคงมีโอกาสที่ไม่เป็นศูนย์ในการทะลุผ่านกำแพงศักย์ไปได้ ปรากฏการณ์นี้เป็นรากฐานของกระบวนการสลายตัวให้อนุภาคแอลฟาในฟิสิกส์นิวเคลียร์ การทำงานของไดโอดทันเนล และกล้องจุลทรรศน์สแกนแบบส่องกราด (Scanning Tunneling Microscope: STM)

\section{ตัวอย่างการคำนวณตามกรอบวิเคราะห์ P-S-C-T}

\begin{psctexample}{6.1}{การคำนวณระดับพลังงานและฟังก์ชันคลื่นของอิเล็กตรอนในบ่อศักย์อนันต์ 1 มิติ}
\textbf{1. ภาพและสถานการณ์ทางกายภาพ (Picture):}\\\\ อิเล็กตรอนถูกขังอยู่ในกล่องศักย์ 1 มิติความกว้าง $L$ โดยมี $V(x) = 0$ สำหรับ $0 \le x \le L$ และ $V(x) = \infty$ ที่บริเวณภายนอก

\vspace{4pt}
\textbf{2. กลยุทธ์และแบบจำลองทางฟิสิกส์ (Strategy):}\\\\ แก้สมการชเรอดิงเงอร์ในบ่อศักย์ กำหนดเงื่อนไขขอบเขต $\psi(0) = \psi(L) = 0$ และการปรับบรรทัดฐาน (Normalization)

\vspace{4pt}
\textbf{3. ขั้นตอนการคำนวณเชิงตัวเลข (Calculation):}\\\\ ในบริเวณ $0 \le x \le L$ สมการชเรอดิงเงอร์คือ:
\begin{equation}
\frac{d^2\psi}{dx^2} + k^2\psi = 0, \quad k = \frac{\sqrt{2mE}}{\hbar}
\end{equation}
ผลเฉลยทั่วไป $\psi(x) = A\sin(kx) + B\cos(kx)$
จากเงื่อนไขขอบเขต $\psi(0) = 0 \implies B = 0$
และ $\psi(L) = 0 \implies \sin(kL) = 0 \implies k_n L = n\pi \quad (n = 1, 2, 3, \dots)$
ได้ระดับพลังงานที่ถูกควอนไตซ์ (Quantized Energy Levels):
\begin{equation}
E_n = \frac{\hbar^2 k_n^2}{2m} = \frac{n^2 \pi^2 \hbar^2}{2mL^2} = \frac{n^2 h^2}{8mL^2}
\end{equation}
จากการปรับบรรทัดฐาน $\int_{0}^{L} A^2 \sin^2(n\pi x/L) dx = 1 \implies A = \sqrt{2/L}$:
\begin{equation}
\psi_n(x) = \sqrt{\frac{2}{L}} \sin\left(\frac{n\pi x}{L}\right)
\end{equation}

\vspace{4pt}
\textbf{4. การทดสอบความสมเหตุสมผล (Test):}\\\\ ระดับพลังงานมีลักษณะไม่ต่อเนื่องและมีพลังงานสถานะพื้น (Zero-point Energy: $E_1 > 0$) ซึ่งสอดคล้องกับหลักความไม่แน่นอนของไฮเซนเบิร์ก $\Delta x \Delta p \ge \hbar/2$
\end{psctexample}

\section{การเชื่อมโยงสู่การวิจัยฟิสิกส์ร่วมสมัย}

\begin{researchbox}{ข้อมูลเชิงลึกจากงานวิจัย}
ในระบบจุดควอนตัม (Quantum Dots) และผลึกนาโน ระดับพลังงานของเอ็กซิตอนถูกควบคุมด้วยขนาดของกล่องตามสมการ $E_n \propto 1/L^2$ ปรากฏการณ์นี้ทำให้เราสามารถปรับเปลี่ยนความยาวคลื่นแสงที่ดูดกลืนและเปล่งออกมาได้โดยการควบคุมขนาดของอนุภาคนาโนในการทดลองสังเคราะห์ทางวัสดุศาสตร์
\end{researchbox}

\section{การฝึกแก้โจทย์ปัญหาเชิงระบบตามกรอบ C-C-A-F}

\begin{ccafsolution}{การวิเคราะห์สัมประสิทธิ์การทะลุผ่านของอิเล็กตรอนผ่านกำแพงศักย์แบบสี่เหลี่ยม}
\textbf{ขั้นที่ 1: การสร้างมโนทัศน์ (Conceptualize):}\\\\ อิเล็กตรอนพลังงาน $E = 2.0\text{ eV}$ วิ่งเข้าหากำแพงศักย์สี่เหลี่ยมความสูง $V_0 = 5.0\text{ eV}$ และความหนา $L = 0.5\text{ nm}$

\vspace{4pt}
\textbf{ขั้นที่ 2: การจำแนกประเภทโจทย์ (Categorize):}\\\\ การคำนวณสัมประสิทธิ์การทะลุผ่านเชิงควอนตัม (Quantum Transmission Coefficient)

\vspace{4pt}
\textbf{ขั้นที่ 3: การวิเคราะห์เชิงคณิตศาสตร์ (Analyze):}\\\\ คำนวณค่าคงที่การลดทอนคลื่น $\kappa$:
\begin{equation}
\kappa = \frac{\sqrt{2m(V_0 - E)}}{\hbar} = \frac{\sqrt{2(9.109 \times 10^{-31})(3.0 \times 1.602 \times 10^{-19})}}{1.055 \times 10^{-34}} \approx 8.87 \times 10^9\text{ m}^{-1}
\end{equation}
สัมประสิทธิ์การทะลุผ่าน $T$:
\begin{equation}
T \approx e^{-2\kappa L} = \exp\left[-2(8.87 \times 10^9)(0.5 \times 10^{-9})\right] = e^{-8.87} \approx 1.4 \times 10^{-4}
\end{equation}

\vspace{4pt}
\textbf{ขั้นที่ 4: การสรุปและประเมินผล (Finalize):}\\\\ อิเล็กตรอนประมาณ 14 ตัวจากทุกๆ 100,000 ตัว สามารถทะลุผ่านกำแพงศักย์ที่มีความสูงมากกว่าพลังงานของตัวมันเองได้ ซึ่งยืนยันธรรมชาติความเป็นคลื่นของสสาร
\end{ccafsolution}

\section{แบบฝึกหัดทบทวนและโจทย์ท้าทายประจำบท}

\begin{enumerate}[leftmargin=1.5cm, label={\textbf{\thechapter.\arabic*}}, itemsep=8pt]
    \item \textbf{โจทย์เชิงแนวคิด (Conceptual):} จงอธิบายความสำคัญทางกายภาพของกฎและสมการหลักในบทเรียนนี้ พร้อมยกตัวอย่างสถานการณ์จริงในธรรมชาติ
    \item \textbf{โจทย์การประมาณค่าเชิงฟิสิกส์ (Estimation):} จงใช้การวิเคราะห์เชิงมิติและอันดับของขนาด (Order of Magnitude) ในการประมาณค่าปริมาณสำคัญที่เกี่ยวข้อง
    \item \textbf{โจทย์วิจัยขั้นสูง (Research Challenge):} จงเขียนสคริปต์ Python จำลองพฤติกรรมของระบบตามสมการที่ได้ศึกษา พร้อมพลอตกราฟเปรียบเทียบผลลัพธ์เชิงทฤษฎี
\end{enumerate}

